% Workshop-style draft (IEEE TPEL / IEEEtran). Structure only: no results are filled in.
% Every number must come from results/ tables produced by src/experiments; see docs/reference_audit.md.
\documentclass[journal]{IEEEtran}
\usepackage{amsmath,graphicx,booktabs,siunitx}

\title{Small State-Space Models for Transient Magnetic Hysteresis:\\
Accuracy--Size Trade-off and Cross-Material Generalization}
\author{Harshil~Agarwal and Mohammad~Nayeem~Teli%
\thanks{Department of Computer Science, University of Maryland, College Park.}}

\begin{document}
\maketitle

\begin{abstract}
TODO after experiments. State: task (MagNet Challenge 2, $H'(t)=f(B,H,B',T)$), model (S4D),
the two questions, the headline numbers with their metric and test set, and the caveats.
\end{abstract}

\section{Introduction}
% Core loss, why transients, MagNet Challenge 2, gap: SSMs barely used by the 32 teams.

\section{Task and Evaluation}
% 16 MHz sequences, 1000-step segments, variable warm-up read per sequence.
% Two metric families (docs/reference_audit.md section 3):
%   tutorial-3: RMSE / mean|H|, energy NRMSE over full span;
%   grading:    RMS(err) / RMS(H), per-sequence energy error over predicted segment / core loss.
% Headline = 95th percentile. Leaderboard comparisons only with the grading family on A--E.

\section{Model}
% Dual input window (B,H) -> pointwise lift -> L x S4D blocks (diagonal SSM, ZOH, S4D-Lin init)
% -> last-step features ++ [B_next, T] -> MLP. Parameter count formula vs d_model, n_layers.
% Training: teacher-forced next-step MSE; selection by rollout p95 on validation segments.

\section{Reproduction of the Baseline}
% Gate: tutorial-3 model (2,399 params) reproduced exactly (33.27 / 16.13 / 39.78).
% Retrained 771-param dual LSTM on our split as the small-model reference.

\section{RQ1: Accuracy vs. Model Size}
% Figure: results/size_sweep/error_vs_params_<M>.pdf. Table: size_sweep_summary.csv.

\section{RQ2: Generalization to Unseen Materials}
% Figure: results/cross_material/cross_material_seq_rms_p95.pdf.
% Zero-shot vs few-shot fine-tune vs in-domain.

\section{Autoregressive Drift}
% Figure: drift curves; error vs steps after warm-up, by warm-up length.

\section{Limitations}
% Subset of materials; own held-out split vs leaderboard A--E; single architecture.

\section{Conclusion}

\bibliographystyle{IEEEtran}
\begin{thebibliography}{9}
\bibitem{mc}M. Chen et al., ``MagNet Challenge for data-driven power magnetics modeling,''
\emph{IEEE Open J. Power Electron.}, 2024, doi: 10.1109/OJPEL.2024.3469916.
\bibitem{magnetx}S. Wang et al., ``MagNetX: Foundation neural network models for simulating power
magnetics in transient,'' in \emph{Proc. IEEE APEC}, 2025, pp. 2438--2445.
\bibitem{magnetxdb}H. Kwon et al., ``MagNetX: Extending the MagNet database for modeling power
magnetics in transient,'' in \emph{Proc. IEEE APEC}, 2025, pp. 566--572.
\bibitem{s4d}A. Gu, A. Gupta, K. Goel, and C. R\'e, ``On the parameterization and initialization of
diagonal state space models,'' in \emph{Proc. NeurIPS}, 2022.
\bibitem{s4}A. Gu, K. Goel, and C. R\'e, ``Efficiently modeling long sequences with structured state
spaces,'' in \emph{Proc. ICLR}, 2022.
\end{thebibliography}
\end{document}
